\section{Boolean-valued construction and the consistency bridge}
\label{sec:annex-p06-boolean-forcing-consistency}

\begin{annexitem}{annex:p06.bfc.completion}{thm:p06-bfc-completion}{Regular-open operations and the dense condition map}
\begin{proof}
For downward open \(U\), the set \(j(U)\) is downward open, includes
\(U\), and is monotone in \(U\). If \(j(U)\) is dense below \(p\),
below any \(q\leq_P p\) first enter \(j(U)\), then enter \(U\).
Thus \(j(j(U))=j(U)\). The set \(n(U)\) is open and regular:
if a condition had an extension in \(U\), the cone below that extension
would have no member of \(n(U)\), so \(n(U)\) could not be dense
below the original condition. Also \(U\cap n(U)=\varnothing\), and
\(U\cup n(U)\) is dense: below a condition either find an extension
in \(U\) or that condition itself belongs to \(n(U)\).

An intersection of any set family of regular open sets is open and
regular, since its regularization is included in every member of the
family. It is their greatest lower bound. The regularization of their
union contains every member and is included in any regular open upper
bound; it is their least upper bound. The empty meet and join are
\(P\) and \(\varnothing\). For regular \(U,V,W\),
\[
 U\cap j(V\cup W)=j((U\cap V)\cup(U\cap W)).
\]
For the forward inclusion, openness of \(U\) keeps the entire cone
inside \(U\), and density of \(V\cup W\) gives the right-hand density.
For the reverse inclusion, monotonicity gives membership in both
\(U=j(U)\) and \(j(V\cup W)\). This proves distributivity; the
complement calculation above proves the Boolean laws. The algebra is
a set by Separation on \(\mathcal P(P)\) and is nontrivial since \(P\)
is nonempty. No topological or class-sized completion has been presumed.

If \(U\in B\) is nonempty, choose \(p\in U\). Its cone lies in \(U\),
so \(0_B<i(p)\leq_B U\). If \(p,q\) have a common strengthening
\(r\), then \(i(r)\leq_B i(p)\cap i(q)\) is nonzero. Conversely,
from a condition in both images first strengthen into \(\downarrow p\),
then into \(\downarrow q\); their defining densities ensure a common
strengthening. For an antichain of nonzero Boolean elements, Choice
selects a condition image below each. Distinct selected conditions are
incompatible by the preceding calculation, so ccc of \(P\) makes the
original Boolean antichain countable. Every operation and quantifier in
this argument concerns sets, so the resulting ZFC theorem applies
internally in a possibly ill-founded model too.
\end{proof}
\end{annexitem}

\begin{annexitem}{annex:p06.bfc.values}{thm:p06-bfc-values}{The bounded tables, all equality cases and Boolean logical soundness}
\begin{proof}
\paragraph{Internal set construction.}
This proof first takes place in ZFC, with no membership model or generic.
Transitive closures are built by an \(\omega\)-sequence and Union.
Rank approximations on transitive subsets satisfy
\(r(x)=\sup_{y\in x}(r(y)+1)\). Two approximations agree: Foundation
on their set of disagreement points gives a point with identical earlier
values. Foundation induction constructs an approximation on each
\(\operatorname{tc}(\{x\})\) by collecting the unique approximations
for its members with Replacement, taking their compatible union, then
adjoining the value at \(x\). Replacement on the whole transitive closure
collects the graphs. This gives unique physical ranks without an assumed
external well-foundedness theorem about a model.

Close a proposed function-style name under the first-coordinate entries
of its function graph. This closure is a set bounded by its transitive
closure. Test that every object reached is a function with coefficients
in \(B\). For accepted names, a proper subname has smaller physical rank.
Ordinal tables on a bounded set of such names construct intrinsic ranks
\(\rho(u)=\sup_{t\in\operatorname{dom}(u)}(\rho(t)+1)\): collect the
newly assigned values with Replacement, and unite earlier tables at
limits. The complete ordinal recursion theorem
\cref{thm:p06-definable-ordinal-recursion} supplies these set tables.
The same theorem constructs the hierarchy of functions on earlier-name
sets; each function graph is bounded by a power set of their product
with \(B\). Thus the name test and hierarchy membership agree by rank
induction and give a first-order predicate.

On a rank-bounded set of names, construct the two atomic value tables in
lexicographic order of the sorted rank pair (maximum, minimum).
The outer ordinal table visits the maximum, and the inner table visits
the minimum. Equality uses membership at pairs with a proper subname;
membership uses equality at such pairs. Each call lowers the measure.
Completeness of \(B\) gives a unique join or meet of each set of
already assigned values. Replacement forms the next table of outputs.
Empty names give the base values: empty membership is zero and empty
equality is one. Induction proves uniqueness and agreement between
larger bounds. A formula asserting existence of a bound and the specified
tables therefore defines the atomic values. These are actual set-bounded
recursions, not an external evaluation of internal names.

For each fixed standard formula, finite recursion through its syntax
gives the displayed logical clauses. At an existential, Separation on
\(B\) collects the realized matrix values before taking their join.
Fresh renaming prevents capture. Boolean complement and the classical
expansions give universal quantification as the meet of its set of
matrix values. This is a separate defining formula for each fixed
formula; there is no quantifier over all formula codes in this construction.

In any complete Boolean algebra finite meets distribute over set joins.
Indeed, if \(z\) bounds every \(c\wedge_B s\), then every \(s\) is
below \(\neg_B c\vee_B z\). Their join is below that element, so
\(c\wedge_B\bigvee S\leq_B z\). The converse upper-bound direction
is immediate. This proves the distribution used in the calculations below.

\paragraph{The joint rank induction for equality.}
Reflexivity follows by rank induction: the membership term at a child
\(t\) includes \(u(t)\wedge_B E_B(t,t)=u(t)\), making every
reflexive equality factor one. Symmetry is in the definition.
Prove transitivity and membership substitution jointly by induction on
the descendingly sorted triple of name ranks. For substitution in the
first membership argument, distribute to obtain
\[
 E_B(u,v)\wedge_B I_B(u,w)
 =\bigvee_{t\in\operatorname{dom}(w)}
 w(t)\wedge_B E_B(u,v)\wedge_B E_B(u,t)
 \leq_B I_B(v,w).
\]
The inequality is lower-triple transitivity at \((v,u,t)\).
For substitution in the second argument, expand \(I_B(u,v)\).
The equality definition gives
\(E_B(v,w)\wedge_B v(t)\leq_B I_B(t,w)\), and lower-triple
first-argument substitution at \((t,u,w)\) gives
\(E_B(u,t)\wedge_B I_B(t,w)\leq_B I_B(u,w)\).
Taking the join proves the second-argument law.

For transitivity, put \(c=E_B(u,v)\wedge_B E_B(v,w)\).
For each \(t\in\operatorname{dom}(u)\),
\(c\wedge_B u(t)\leq_B E_B(v,w)\wedge_B I_B(t,v)
\leq_B I_B(t,w)\), using the lower-triple second-argument law.
Interchange \(u,w\) for the other inclusion. Both equality factors
for \(E_B(u,w)\) therefore dominate \(c\). In each displayed use
one original name was replaced by a proper subname; sorting strictly
decreases the triple. Empty domains impose no factors. This discharges
all cases of the joint induction.

With \(c\leq_B E_B(u,v)\), these laws give equal membership/equality
values after meeting with \(c\). Induction on a fixed formula extends
this to all formulas. For negation use relative Boolean complement below
\(c\); for conjunction use finite meets; for existential quantification
use the just-proved distribution over its set of realized values. This
also proves the bounded-quantifier calculation
\[
 \lVert\exists x\in u\ \theta(x)\rVert_B
 =\bigvee_{t\in\operatorname{dom}(u)}u(t)\wedge_B\lVert\theta(t)\rVert_B.
\]
Expand membership in \(u\), substitute equals under each summand,
then take joins; the reverse inequality uses the instance \(x=t\)
and reflexive equality. The universal version follows by complement.

\paragraph{Every logical inference.}
For a derivation under open assumptions let \(g\) be their finite meet
of values. An assumption leaf has \(g\) below its value. Conjunction
introduction/elimination, disjunction introduction, implication elimination
and negation elimination are the respective Boolean inequalities.
Disjunction elimination uses distribution of \(g\) over the two alternatives.
Implication or negation introduction removes the discharged assumption
using \(g\wedge_B a\leq_B c\) iff \(g\leq_B\neg_B a\vee_B c\).
Bottom elimination is \(0_B\leq_B c\); classical reductio uses
\(g\wedge_B\neg_B c=0_B\) iff \(g\leq_B c\).
Universal elimination and existential introduction are meet/join bounds.
For universal introduction, freshness makes \(g\) independent of the
generalized variable; if \(g\) is below every instance, it is below
their meet. For existential elimination, freshness makes both \(g\)
and the conclusion independent of the eigenvariable; the minor derivation
gives \(g\wedge_B\lVert\theta(t)\rVert_B\leq_B c\) for each name,
and distribution gives the same inequality for the existential join.
Equality introduction is reflexivity; equality elimination is substitution.
Capture-avoiding alpha conversion changes no value by the formula induction.
Repeated assumptions use idempotence of meet; the actual calculus's stronger
nonvacuous-discharge and freshness conditions only restrict these valid cases.
Thus every constructor of its classical natural-deduction proof grammar
preserves the required inequality. This proof is internal Boolean algebra,
not soundness obtained from a presumed quotient model.
\end{proof}
\end{annexitem}

\begin{annexitem}{annex:p06.bfc.fullness}{thm:p06-bfc-fullness}{The set-bound maximum principle}
\begin{proof}
For an antichain \(A\) and names \(u_a\), put
\(D=\bigcup_{a\in A}\operatorname{dom}(u_a)\) and define the mixed
function \(u\) on \(D\) by
\[
 u(t)=\bigvee_{a\in A}a\wedge_B u_a(t),
\]
with zero for absent coefficients. Replacement, Union and completeness
form its graph; the collected ranks of the names bound its subnames.
For fixed \(a\in A\), disjointness gives
\(a\wedge_B u(t)=a\wedge_B u_a(t)\). Since
\(u_a(t)\leq_B I_B(t,u_a)\) and \(u(t)\leq_B I_B(t,u)\), both
atomic equality inclusion tests dominate \(a\); thus \(a\leq_B E_B(u,u_a)\).
The empty family yields the empty name.

For \(c=\lVert\exists x\theta(x,\bar u)\rVert_B\), let
\(S\subseteq B\) be its set of realized matrix values. The positive
elements below members of \(S\) are dense below \(c\): a nonzero
\(d\leq_B c\) meets some member of \(S\) positively, by finite-meet
distribution. Well-order this dense set by ground Choice and retain
an element just when it is disjoint from all previously retained ones.
Ordinal recursion and Separation give an antichain \(A\). If
\(c\wedge_B\neg_B\bigvee A\) were nonzero, density would give an
element below it which the construction must have retained, a contradiction.
Therefore \(\bigvee A=c\).

For each \(a\in A\), a witnessing name exists with
\(a\leq_B\lVert\theta(u_a,\bar u)\rVert_B\). This is a set-indexed
choice from an initially class-sized pool. To bound it, assign to each
\(a\) the least ordinal stage containing a witness. Its leastness is
a functional first-order formula, so Replacement collects the stages.
Their supremum plus one bounds a set of candidate names. Choice on this
bounded family supplies the \(u_a\). This is also a direct proof of the
ground Collection instance being used. Mix the family. Formula congruence
gives \(a\leq_B\lVert\theta(u,\bar u)\rVert_B\) for each \(a\);
their join gives \(c\leq_B\lVert\theta(u,\bar u)\rVert_B\).
The reverse inequality is the definition of \(c\) as the join of all
realized values. If \(c=0_B\), every matrix value is zero and the empty
name is already a witness. No filter or quotient has been used.
\end{proof}
\end{annexitem}

\begin{annexitem}{annex:p06.bfc.axioms}{thm:p06-bfc-axioms}{Every internal axiom-value calculation}
\begin{proof}
We work in ZFC on Boolean values, before taking any quotient.
Write \(v(\varphi)=\lVert\varphi\rVert_B\).
Reflexivity and substitution show that
\[
 v(\forall x(x\in u\leftrightarrow x\in w))=E_B(u,w).
\]
The forward inequality follows by testing each child of \(u,w\), whose
membership value dominates its coefficient; the reverse is membership
congruence. Thus Extensionality has value one. Empty Set uses the empty
function, whose membership values are all zero.

The name \(\{u,w\}\) with coefficient one at each of its one or two
literal domain names has membership value \(E_B(x,u)\vee_B E_B(x,w)\).
This proves Pairing. Applying this constructor twice supplies the
Kuratowski ordered-pair name \(O(u,w)\). For Union of \(u\), use the
set \(D=\bigcup_{t\in\operatorname{dom}(u)}\operatorname{dom}(t)\)
and coefficient at \(s\in D\)
\[
 \bigvee_{t\in\operatorname{dom}(u)}u(t)\wedge_B t(s).
\]
Its membership value at \(x\) is
\(\bigvee_t u(t)\wedge_B I_B(x,t)\), precisely the bounded existential
\(v(\exists t\in u\ x\in t)\). The set bounds and every coefficient
join are explicit, so the result is a name, not an unbounded union expression.

For Separation on \(u\) at a fixed formula \(\theta\), take the
function on \(\operatorname{dom}(u)\) with coefficient
\(u(t)\wedge_B v(\theta(t))\). Distribution and formula congruence
give its membership value at every \(x\) exactly
\(I_B(x,u)\wedge_B v(\theta(x))\). This proves every parameterized
Separation instance.

For Power Set, the functions with domain a subset of
\(\operatorname{dom}(u)\) and coefficients in \(B\) form a set
\(C\subseteq\mathcal P(\operatorname{dom}(u)\times B)\). They are
all names. Give each \(w\in C\) coefficient \(v(w\subseteq u)\)
in a new name \(p\). Congruence gives \(I_B(z,p)\leq_B v(z\subseteq u)\).
For the reverse inequality form the particular \(z'\in C\) with
coefficients \(u(t)\wedge_B I_B(t,z)\). Its membership values are
\(I_B(x,u)\wedge_B I_B(x,z)\), by the Separation calculation.
Consequently \(v(z\subseteq u)\leq_B E_B(z,z')\) by the already
proved extensionality identity, and \(v(z'\subseteq u)=1_B\).
The \(z'\) summand of \(I_B(z,p)\) therefore dominates
\(v(z\subseteq u)\). These two inequalities prove the exact Power Set
axiom, not merely a set containing its members.

Canonical names are constructed by the same bounded rank tables with
coefficient one at each checked member. Joint rank induction gives
\(E_B(\check x,\check y)=1_B\) or \(0_B\) according as \(x=y\)
or not, and likewise for membership. Bounded-quantifier reduction then
proves the analogous statement for every fixed bounded ground formula.
The checked \(\omega\) is inductive: checking commutes with the
empty-set, pair and union constructors, since their displayed membership
values agree, and the ground successor-closure instances have value one.
Hence Infinity has value one. The bounded ordinal predicate likewise
makes each checked ordinal an ordinal with value one.

For Foundation, put \(a=v(u\ne\varnothing)\). If a nonzero
\(b\leq_B a\) were disjoint from the value of the required minimal-member
existential, choose a name \(t\) of least intrinsic rank such that
\(b\wedge_B I_B(t,u)>0_B\). Such a name exists, since
\(a=\bigvee_{s\in\operatorname{dom}(u)}u(s)\); first take a positive
summand, then minimize within the ordinal bound of that initial witness.
Every \(s\in\operatorname{dom}(t)\) has lower rank, so
\(b\wedge_B I_B(s,u)=0_B\). The bounded-quantifier calculation yields
\(b\wedge_B v(\exists s\in t\ s\in u)=0_B\).
Thus the nonzero \(c=b\wedge_B I_B(t,u)\) is below the minimal-member
existential value, a contradiction. This proves the Foundation implication
at every name, using internal ordinal minimization, not external evaluation.

For Collection, fix \(\theta(x,y)\). Fullness supplies for each
\(t\in\operatorname{dom}(u)\) a name \(w_t\) realizing the value
\(v(\exists y\theta(t,y))\). The least-stage Collection bound and
ground Choice used in the preceding proof give this as a set family.
Let \(w\) contain these names with coefficient one. If
\(a=v(\forall x\in u\exists y\theta(x,y))\), then for every child
\(t\), \(a\wedge_B u(t)\leq_B v(\theta(t,w_t))\), and
\(I_B(w_t,w)=1_B\). Bounded-quantifier reduction therefore proves
\(a\leq_B v(\forall x\in u\exists y\in w\theta(x,y))\).
For Replacement with uniqueness, the same bound suffices for every image:
on the Boolean part where \(x\) is equal to a child \(t\), uniqueness
identifies any \(\theta(x,y)\)-image with \(w_t\). Formally, multiply
\(a\wedge_B u(t)\wedge_B E_B(x,t)\wedge_B v(\theta(x,y))\)
by the uniqueness implication; it is below \(E_B(y,w_t)\), hence below
\(I_B(y,w)\). Join over the children. Separation on \(w\) with
\(\exists x\in u\theta(x,y)\) now produces the exact image, conditional
on \(a\), so the functional Replacement implication has value one.

Finally well-order \(\operatorname{dom}(u)\) in the ground, enumerating
it as \((t_\gamma)_{\gamma<\alpha}\). The name with coefficient one
at \(O(\check\gamma,t_\gamma)\) is a function on \(\check\alpha\)
whose range contains \(u\), with value one. To verify this, the
ordered-pair identity is available by Boolean logical soundness from
the Extensionality/Pairing values already proved: its ordinary complete
Kuratowski proof uses just those axioms. At a fixed checked index
\(\gamma\), pair membership with output \(y\) reduces to
\(E_B(y,t_\gamma)\), since distinct checked indices have equality value
zero. Existence follows with \(y=t_\gamma\), and equality transitivity
gives uniqueness. Any range point of \(u\) is covered because
\(I_B(x,u)\leq_B\bigvee_\gamma E_B(x,t_\gamma)\).
The now verified ZF axioms prove that a set covered by an ordinal-indexed
function can be well-ordered: take its least-preimage injection by Separation
and Replacement, and pull back the ordinal order. They then prove Choice
by selecting the least member of each factor under a well-order of their
union. Boolean logical soundness transports these actual derivations;
every input condition has value one by the constructions just given.
Thus the Choice axiom also has value one. No axiom has been inferred
from the mere existence of a class of names or from an assumed quotient.
\end{proof}
\end{annexitem}

\begin{annexitem}{annex:p06.bfc.quotient}{thm:p06-bfc-quotient}{An external satisfaction structure from internal names}
\begin{proof}
Interpret the preceding fixed-formula ZFC constructions in \(\mathcal M\).
The interpreted Boolean operations satisfy their ordinary finite algebraic
laws on the actual set \(B^{\mathcal M}\). A join computed internally
for an encoded set is a least upper bound for its interpreted members:
the internal least-upper-bound formula quantifies over exactly those members
and all actual elements of \(B^{\mathcal M}\). No completeness for
other external subsets is asserted. Fullness, interpreted at each fixed
standard formula, supplies an actual witness in \(\mathcal D_b\).

Equality reflexivity/symmetry/transitivity and closure of \(U\) under
finite meets make \(\sim_U\) an equivalence relation. Membership congruence
makes \(\in_U\) independent of representatives. Its classes are ordinary
subsets of the external set \(\mathcal D_b\), so their collection and
membership relation form an actual set-sized structure, even if the
internal name-membership relation is ill-founded externally.
Induct on a standard formula. Atoms hold by definition. Negation uses
the fact that exactly one of a Boolean value and its complement belongs
to \(U\); conjunction uses finite-meet closure and upward closure.
For an existential whose value belongs to \(U\), interpreted fullness
gives a name attaining that same value, hence a witness class by induction.
Conversely a witness class has its matrix value in \(U\), and that value
is below the existential value, so upward closure supplies the latter.
The fixed expansions handle the other constructors. This establishes actual
satisfaction, not a declaration that a quotient should satisfy a theory.
For each standard ZFC axiom separately, \(\mathcal M\) satisfies the
first-order axiom-value theorem just proved; its value is the unit in \(U\).
Thus the quotient satisfies every axiom. Internal Foundation was used only
to prove internal construction statements, never to collapse the model or
to evaluate its names externally.
\end{proof}
\end{annexitem}

\begin{annexitem}{annex:p06.bfc.cohen}{thm:p06-bfc-cohen}{The cardinal argument entirely in Boolean values}
\begin{proof}
The complete Delta-system and Cohen-ccc proofs
\cref{thm:p06-delta-system,thm:p06-add-ccc} are ZFC set arguments, so
the completion theorem makes \(B\) ccc internally. Write \(v\) for
Boolean value. We first give the preservation argument rather than importing
one from an externally evaluated extension.

Let \(\mu\) be an uncountable ground cardinal and \(\lambda<\mu\).
Suppose \(b=v(f:\check\lambda\twoheadrightarrow\check\mu)>0_B\)
for some name \(f\). For \(\alpha<\lambda,\beta<\mu\), put
\(b_{\alpha\beta}=b\wedge_B v(f(\check\alpha)=\check\beta)\),
where graph/function notation has its usual membership-language expansion.
At each \(\alpha\), the nonzero values form an antichain: functionality
and equality of distinct checked ordinals give disjointness. Their join
is \(b\), by totality and the bounded-quantifier computation for
\(\check\mu\). Thus
\(S_\alpha=\{\beta:b_{\alpha\beta}>0_B\}\) is countable by ccc.
Ground Choice chooses injections of these sets into \(\omega\).
Mapping a member of their union to its least occurring index and its
chosen natural index injects that union into \(\lambda\times\omega\).
The ground infinite-cardinal finite-coding theorem
\cref{thm:infinite-cardinal-finite-coding} gives its cardinality
at most \(\max(|\lambda|,\aleph_0)<\mu\).
Choose \(\beta_*\) outside this union. At every index its function-value
Boolean element meets \(b\) in zero; surjectivity would make their join
cover \(b\), a contradiction. Hence all these surjection values are zero.
Bounded existential reduction over \(\check\mu\) shows that no name
for a smaller ordinal can supply a counterexample: it reduces to the
join of the cases \(\lambda<\mu\) just proved. Therefore each checked
uncountable ground cardinal remains a cardinal with value one.

For \(\mu=\aleph_1\), every ground \(\lambda<\mu\) has an injection
into \(\omega\); its checked graph witnesses this with value one by
bounded ground-formula preservation. Bounded universal reduction makes
the assertion true for every ordinal below \(\check\mu\), not only
for individually named ground indices in a quotient. The checked
\(\omega\) is the least infinite ordinal: its bounded ordinal,
successor-closure and nonzero-member-successor properties are preserved,
and the ZF ordinal argument characterizes \(\omega\) from these properties.
Thus \(v(\check\aleph_1=\aleph_1)=1_B\).
Apply the same argument with injections of ground ordinals below
\(\aleph_2\) into \(\aleph_1\); preservation of \(\aleph_2\) and
the absence of an intervening cardinal give
\(v(\check\aleph_2=\aleph_2)=1_B\).

For each \(\alpha<\kappa\) define a name \(r_\alpha\) on the checked
naturals with coefficient
\(i(\{((\alpha,n),1)\})\) at \(\check n\). Its complement is the
image of the opposite bit condition: these two images are disjoint and
their union is dense. Thus \(v(r_\alpha\subseteq\check\omega)=1_B\).
For distinct \(\alpha,\beta\), below any finite condition choose a fresh
\(n\) and give its two coordinates opposite bits. The image of that
extension is below the Boolean value of a difference between the two
names. Density of condition images implies that the join of all such
difference values is one; the atomic equality clauses therefore give
\(E_B(r_\alpha,r_\beta)=0_B\).
The ground set of ordered-pair names \(O(\check\alpha,r_\alpha)\),
all with coefficient one, is a name for an injection from
\(\check\kappa\) into the subsets of \(\check\omega\): at each
checked index pair membership reduces to equality with its assigned
output, giving functionality, and the preceding zero equalities give
injectivity. The same bounded-quantifier calculation extends the statement
to all index names. Ordered-pair identity follows from the Boolean-valued
Extensionality/Pairing argument used in the axiom proof.
Since the checked two successor cardinals are the two actual successor
cardinals in the Boolean theory, this injection contradicts CH there.
Boolean logical soundness gives \(v(\neg\mathsf{CH})=1_B\).
All arguments were ZFC calculations of values, with no generic filter
or external well-founded evaluation.
\end{proof}
\end{annexitem}

\begin{annexitem}{annex:p06.bfc.consistency}{thm:p06-bfc-consistency}{Completeness, a proper ultrafilter and the target model}
\begin{proof}
Assume syntactic consistency of ZFC. The existing countable Henkin
completeness construction \cref{cor:fol-completeness-countable}, applied
to the countable membership language and ZFC theory, gives an actual
countable membership structure \(\mathcal M=(D,E)\models\mathsf{ZFC}\).
No Mostowski collapse or transitive-model assertion is made.
Inside this structure take its Cohen poset and regular-open completion.
The preceding ZFC theorems give a nontrivial interpreted Boolean algebra,
the interpreted axiom values, fullness, and the non-CH value one.

Enumerate its actual Boolean elements as \(a_0,a_1,\ldots\). Set
\(c_0=1_B\). At stage \(n\), if \(c_n\wedge_B a_n\ne0_B\), use
that meet for \(c_{n+1}\); otherwise use
\(c_n\wedge_B\neg_B a_n=c_n\). Each \(c_n\) is nonzero and the
sequence descends. The set
\(U=\{a:\exists n\ c_n\leq_B a\}\) is a proper filter, and the
construction decides each element or its complement, so it is an
ultrafilter. It is not required to be an internally encoded set or a generic.
The actual quotient structure of \cref{thm:p06-bfc-quotient} satisfies
all standard ZFC axioms and \(\neg\mathsf{CH}\), whose value is the unit
in \(U\). First-order soundness rules out a contradiction derivation
from that theory. This proves the bare syntactic consistency implication
from its original premise, with an explicit satisfaction structure.
It does not claim a primitive-recursive arithmetized PA proof translation.
\end{proof}
\end{annexitem}
